% uid: ""
% school: "lycee"
% level: "3eme"
% section: ["Economie et Gestion"]
% chapter: "Suites réelles"
% title: "Piège : raison négative, pas de limite"
% difficulty: "Difficile"
% points: null
% tags: ["suites", "suite-recurrente", "raison-negative", "absence-de-limite", "piege"]
% provenance: "own"
% dressing: "U(n+1)=aU(n)+b, a<-1"
% rung: 4
% batch: "2026-11-20-suites-3eco-A"

% ---STATEMENT---
Soit $(U_n)$ la suite définie par $U_0=1$ et $U_{n+1}=-1{,}5\,U_n+5$.

\begin{enumerate}
  \item Calculer $U_1$ et $U_2$.
  \item Une élève affirme : « si $(U_n)$ a une limite $\ell$, alors $\ell=-1{,}5\,\ell+5$, donc $\ell=2$. » Que vaut ce raisonnement ?
  \item Trouver $\alpha$ tel que $\alpha=-1{,}5\,\alpha+5$ et poser $V_n=U_n-\alpha$. Montrer que $(V_n)$ est géométrique et exprimer $U_n$.
  \item La suite $(U_n)$ a-t-elle une limite ? Conclure sur l'affirmation de l'élève.
\end{enumerate}

% ---SOLUTION---
\begin{enumerate}
  \item $U_1=-1{,}5\times1+5=3{,}5$ et $U_2=-1{,}5\times3{,}5+5=-5{,}25+5=-0{,}25$.
  \item Si $(U_n)$ converge vers $\ell$, alors $U_{n+1}\to\ell$ aussi et on obtient $\ell=-1{,}5\,\ell+5$, soit $\ell=2$. C'est donc le \emph{seul candidat possible}.

  Mais ce raisonnement \textbf{suppose} que la limite existe : il ne le prouve pas.
  \item $\alpha=-1{,}5\,\alpha+5\iff2{,}5\,\alpha=5\iff\alpha=2$.
  $$V_{n+1}=U_{n+1}-2=-1{,}5\,U_n+3=-1{,}5\,(U_n-2)=-1{,}5\,V_n.$$
  $(V_n)$ est géométrique de raison $-1{,}5$ avec $V_0=1-2=-1$. Donc $V_n=-(-1{,}5)^n$ et
  $$ \ib{$U_n=2-(-1{,}5)^n$}.$$
  \emph{Vérification} : $U_1=2+1{,}5=3{,}5$ $\checkmark$ \quad $U_2=2-2{,}25=-0{,}25$ $\checkmark$
  \item La raison $-1{,}5\le-1$ : $(-1{,}5)^n$ vaut $1;\,-1{,}5;\,2{,}25;\,-3{,}375;\dots$ Les signes alternent et les valeurs absolues tendent vers $+\infty$. Donc $(U_n)$ \textbf{n'a pas de limite}.

  L'élève avait trouvé le bon \emph{candidat} ($2$ est le point d'équilibre : si $U_0=2$ alors $U_n=2$ pour tout $n$), mais la suite ne s'en approche pas, elle s'en éloigne en oscillant.

  \textbf{Piège :} résoudre $\ell=a\ell+b$ ne démontre jamais l'existence de la limite. Il faut étudier la raison $a$ : si $|a|<1$ la limite existe et vaut $\alpha$ ; si $a\le-1$ elle n'existe pas (sauf $U_0=\alpha$).
\end{enumerate}

\textbf{Idée :} $\alpha$ n'est qu'un candidat : c'est la raison $a$ qui décide si la suite s'en approche.
